\renewcommand{\bibname}{Quellen und Bandzuordnung}
\begin{thebibliography}{9}
\addcontentsline{toc}{chapter}{Quellen und Bandzuordnung}
\footnotesize
\setlength{\itemsep}{5pt}
\setlength{\parskip}{0pt}
\interlinepenalty=10000

\bibitem{project}
Martin Kunze: \emph{Die Grundlagen der Mathematik}.
Band~13: untere und obere Schrankenmengen; Infimum und Supremum;
Charakterisierung eines Kandidaten und der Gleichheit mit einem
vorhandenen Wert; Dualität der Schranken- und Grenzbegriffe.
Die getrennten Einführungs- und Auswertungsregeln des Hauptbands
sind hier zu Charakterisierungen zusammengefaßt.
Band~12 liefert die vorausgesetzte Halbordnung.
Format nach den Leitfäden zu den Bänden~3, 4, 5, 28 und~37.

\bibitem{defs}
The mathlib Community: \emph{Mathlib.Order.Bounds.Defs},
offizielle Dokumentation, abgerufen am 23.~Oktober 2026.
Begriffe \path{upperBounds}, \path{lowerBounds}, \path{BddAbove},
\path{BddBelow}, \path{IsLUB} und \path{IsGLB} zum Abgleich von
Kapitel~\ref{b13guide:ch:foundations}.
\url{https://leanprover-community.github.io/mathlib4_docs/Mathlib/Order/Bounds/Defs.html}.

\bibitem{basic}
The mathlib Community: \emph{Mathlib.Order.Bounds.Basic},
offizielle Dokumentation, abgerufen am 23.~Oktober 2026.
Vergleichsstellen für Kapitel~\ref{b13guide:ch:bounds} bis
\ref{b13guide:ch:iteration}: insbesondere
\path{upperBounds_union}, \path{upperBounds_empty},
\path{isLUB_iff_le_iff}, \path{IsLUB.upperBounds_eq},
\path{IsLUB.unique}, \path{isGreatest_iff_isLUB_and_mem},
\path{isLUB_empty_iff}, \path{IsLUB.of_isCofinalFor},
\path{isGLB_le_isLUB}, \path{isLUB_lowerBounds} und
\path{isGLB_upperBounds} sowie die jeweiligen dualen Aussagen.
Für die ergänzten Intervall- und strikten Vergleichsregeln:
\path{bddBelow_bddAbove_iff_subset_Icc}, \path{upperBounds_Icc},
\path{lowerBounds_Icc}, \path{isLUB_Icc}, \path{isGLB_Icc},
\path{isLUB_lt_iff} und \path{lt_isGLB_iff}.
\url{https://leanprover-community.github.io/mathlib4_docs/Mathlib/Order/Bounds/Basic.html}.

\bibitem{families}
The mathlib Community: \emph{Mathlib.Order.Bounds.Lattice},
offizielle Dokumentation, abgerufen am 23.~Oktober 2026.
Einträge \path{upperBounds_iUnion}, \path{lowerBounds_iUnion},
\path{isLUB_iUnion_iff_of_isLUB} und
\path{isGLB_iUnion_iff_of_isGLB} zum Abgleich von
Theorem~\ref{b13guide:thm:family-bounds} und
Kapitel~\ref{b13guide:ch:families}.
Diese Einträge benötigen trotz des Modulnamens keine
Verbandaxiome für den Träger.
\url{https://leanprover-community.github.io/mathlib4_docs/Mathlib/Order/Bounds/Lattice.html}.


\newpage
\bibitem{wodzicki}
Mariusz Wodzicki: \emph{Notes on Ordered Sets},
University of California, Berkeley, 30.~September 2018.
Abschnitte~1.3--1.4: Schrankenmengen und allgemeine Infima/Suprema;
Aufgabe~13 und Abschnitt~1.4.3: Sandwichform in
Theorem~\ref{b13guide:thm:sandwich};
Abschnitt~1.8.6: gegenseitige Überdeckung in
Theorem~\ref{b13guide:thm:domination}.
Abschnitt~1.9.3: Abhängigkeit der Schrankenmengen vom Träger.
\url{https://math.berkeley.edu/~wodzicki/1.pdf}.

\bibitem{mizar}
Grzegorz Bancerek: \emph{Directed Sets, Nets, Ideals, Filters, and Maps}.
Mizar Mathematical Library, Artikel \path{WAYBEL_0},
eingereicht am 12.~September 1996; aktuelle Bibliotheksfassung
abgerufen am 23.~Oktober 2026.
Ausschließlich Sätze~31--33 und~36--38: unveränderte Schranken und
Grenzen beim Hinzufügen unterhalb beziehungsweise oberhalb liegender
Elemente, Theorem~\ref{b13guide:thm:one-sided-enlargement}.
Die weiteren Themen dieses Artikels werden nicht übernommen.
\url{https://mizar.uwb.edu.pl/version/current/html/waybel_0.html}.

\bibitem{minimalbounds}
The mathlib Community: \emph{Mathlib.Order.Minimal},
offizielle Dokumentation, abgerufen am 23.~Oktober 2026.
Eintrag \path{minimal_iff_forall_lt} zum Begriff in
Definition~\ref{b13guide:def:minimal-bound}.
Das Kriterium in Theorem~\ref{b13guide:thm:minimal-bound} ist hier
unmittelbar für obere beziehungsweise untere Schranken formuliert.
\url{https://leanprover-community.github.io/mathlib4_docs/Mathlib/Order/Minimal.html}.

\bibitem{minimax}
Yakov Shklarov: \emph{The Minimax Inequality for Lattices},
2.~Januar 2020, Seite~1: schwache Minimax-Ungleichung.
Vergleichsquelle zu Theorem~\ref{b13guide:thm:minimax}; die Quelle
formuliert für Verbände mit den benötigten Grenzen.
Hier wird ausschließlich die aus den Schrankencharakterisierungen
folgende Fassung in einer beliebigen Halbordnung verwendet,
mit expliziter Existenz sämtlicher genannter innerer und äußerer Werte.
Die weiteren Anwendungen und algebraischen Gesetze der Quelle
sind nicht Bestandteil dieses Leitfadens.
\url{https://yakov.shklarov.com/writing/minimax.pdf}.

\end{thebibliography}

\begingroup
\footnotesize
\noindent\textbf{Einordnung.}
Die Quellen dienen dem Abgleich von Begriffen und zentralen Sätzen.
Die Schrankenidentitäten, exakten Vergleichskriterien,
Zwischenmengenregeln und Aussagen zum Trägerwechsel sowie die
Vertauschungsregel und die lokale Minimax-Ungleichung sind hier
als unmittelbare Folgerungen aus den angegebenen Definitionen
und Charakterisierungen zusammengefaßt; nicht jede Formulierung
steht in den Quellen wörtlich. Die Bibliothek behandelt daneben
stärkere Strukturen. Übernommen werden ausschließlich Aussagen,
die in einer beliebigen Halbordnung mit den jeweils genannten
lokalen Existenzbedingungen gelten. Insbesondere werden keine
Resultate mit zusätzlichen Totalitäts-, Halbverbands- oder
Vollständigkeitsannahmen übertragen. Die Verwendung von $\sup S$
und $\inf S$ setzt hier stets deren Existenz voraus; hiervon
abweichende Konventionen für überall definierte Operationen
werden nicht übernommen. Es werden weder neue Forschungsergebnisse
noch historische Erstentdeckungen beansprucht.
\endgroup
